% !TeX root = ../../Book2.tex
\section{限制算子}
\label{b2:sec:2301}

Fourier 变换在 \(L^2\) 中是一个等价类，而光滑曲面在环境空间中具有零 Lebesgue 测度。因此“把 \(L^2\) Fourier 变换代入曲面”没有天然意义。限制问题问的是：从 Schwartz 函数开始的曲面取值操作，是否能由某个环境 \(L^p\) 范数控制，从而连续延拓。

\subsection{\texorpdfstring{\(\widehat f|_S\)}{Fourier 变换在 S 上的限制} 的意义}
\label{b2:subsec:230101}

对紧致光滑超曲面 \(S\subset\mathbb R^d\)、\(d\ge2\)，先定义
\[
 Rf=\widehat f|_S,\qquad f\in\mathcal S(\mathbb R^d).
\]
\symidx{\(Rf\)：曲面 Fourier 限制}
一个限制估计是
\(\|Rf\|_{L^r(S,\mathrm d\sigma)}\le C\|f\|_{L^p(\mathbb R^d)}\)。
若 \(p<\infty\)，稠密性给出唯一的 \(L^p\to L^r(S)\) 延拓。这个延拓是 \(L^r(S)\) 中的极限，并非对每个输入都先选择一个 Fourier 代表再逐点代入。

\ProseParagraphBreak
不能从 Plancherel 直接得到 \(p=r=2\)。取一个曲面坐标邻域，令
\(\widehat f_\epsilon(\xi)=b(\xi)\chi(\rho(\xi)/\epsilon)\)，其中 \(\rho=0\) 定义曲面，\(b|_S\ne0\)、\(\chi(0)=1\)。
曲面限制范数不随 \(\epsilon\) 变化，而
\(\|f_\epsilon\|_2=\|\widehat f_\epsilon\|_2=O(\epsilon^{1/2})\)。
这排除了 \(L^2\to L^2(S)\) 界，也显示环境等价类在零测集取值的困难具有真实的范数后果。

若额外要求空间矩可积，普通平方迹也可以成立。这个模型有助于把限制定理与传统迹定理区别开来。
\begin{proposition}\label{b2:prop:23-moment-trace}
对 Schwartz 输入 \(f\)，有
\[
 \|\widehat f|_{S^{d-1}}\|_2^2
 \leq C_d\|f\|_2\bigl(\|f\|_2+\||x|f\|_2\bigr).
\]
右侧有限的加权 \(L^2\) 空间因而有连续的球面迹延拓。
\end{proposition}
\begin{proof}
对每个方向 \(\omega\)，令 \(F(r)=\widehat f(r\omega)\)，在 \(1/2<r<3/2\) 上使用微积分基本定理。把起点 \(s\) 在这个区间平均，得到
\[
 |F(1)|^2\leq C\int_{1/2}^{3/2}|F(r)|^2\,\mathrm dr
 +C\int_{1/2}^{3/2}|F(r)|\,|F'(r)|\,\mathrm dr.
\]
对球面方向积分，并注意此半径区间的 \(r^{d-1}\) 与一可比。Cauchy--Schwarz 将第二项控制为
\(C\|\widehat f\|_{L^2(\text{壳})}\|\nabla\widehat f\|_{L^2(\text{壳})}\)。
Plancherel 及 \(\partial_{\xi_j}\widehat f=-2\pi\mathrm i\widehat{x_jf}\) 就给出结论。光滑紧支集逼近在带矩的平方范数中成立，估计保证其迹为 Cauchy，完成延拓。
\end{proof}
这条迹界仍是平方理论，但输入空间额外记录位置衰减或矩。Stein--Tomas 改用更小的普通 Lebesgue 指数，从而不要求输入有空间矩；它用曲面的振荡几何换取迹的控制。把上述带矩界写成无条件 \(L^2\) 限制，就会重新遇到薄壳反例，不能省去右侧的矩范数。

\subsection{曲面测度}
\label{b2:subsec:230102}

若 \(S\) 局部写成 \(\xi_d=h(\xi')\)，则
\[
 \mathrm d\sigma=\sqrt{1+|\nabla h|^2}\,\mathrm d\xi'.
\]
对光滑正密度 \(a\)，也可使用 \(\mathrm d\mu=a\,\mathrm d\sigma\)；在 \(a\) 有正上下界的紧片上，两种范数等价。若密度在某点消失，某些必要条件可能改变，所以陈述锐利估计时应明确测度。

\ProseParagraphBreak
余面积公式与频率薄壳联系紧密。对于球面，极坐标给出
\[
 \int_{\mathbb R^d}F(\xi)\,\mathrm d\xi
   =\int_0^\infty r^{d-1}\int_{S^{d-1}}F(r\omega)\,\mathrm d\sigma(\omega)\,\mathrm dr.
\]
在半径 \(1\) 的薄壳中，\(r^{d-1}\) 是有统一上下界的因子。曲面取值可以视为频率薄壳平均在厚度趋零时的极限，但没有范数控制时，这个极限仍可能发散。

\subsection{曲率}
\label{b2:subsec:230103}

本章首先以单位球面为模型，其 Gaussian 曲率处处非零。图形曲面的曲率条件可在驻点处通过 \(D^2h\) 的行列式表达；\chapref{b2:ch:21}已经解释它怎样产生 \((d-1)/2\) 个衰减幂次。
曲率对限制估计的作用并非单纯保证 Fourier 衰减。小曲面帽的法向方向还形成细长的相干管束；所有帽的叠加会产生比单个核估计复杂得多的几何问题。

\ProseParagraphBreak
平面完全不同：对 \(\xi_d=0\)，限制是输入沿 \(x_d\) 积分后在其余变量上 Fourier 变换。把输入写为 \(g(x')h(x_d/L)\)，曲面限制含一个 \(L\) 因子，而 \(L^p\) 范数只含 \(L^{1/p}\)。让 \(L\to\infty\) 排除任何 \(p>1\) 的统一界（取固定非零 \(g,h\)）。这不是一个曲率不足的小损失，而是沿平面法向可以无限相干。

\subsection{缩放}
\label{b2:subsec:230104}

令 \(S_R=RS\)，则面测度缩放为 \(R^{d-1}\mathrm d\sigma\)。
若单位曲面的 \(L^p\to L^r\) 估计成立，换元得到
\[
 \|\widehat f\|_{L^r(S_R)}\le
 C R^{(d-1)/r-d/p'}\|f\|_p.
\]
证明取 \(g(x)=f(x/R)\)，则 \(\widehat g(\xi)=R^d\widehat f(R\xi)\)，
\(\|g\|_p=R^{d/p}\|f\|_p\)，再代入面测度因子。
这一式子把半径依赖写清楚；单位球面本身不是尺度不变对象，不能把不同半径的常数都无条件设成同一个数。
